% TODO: update the executive summary once flight test results are finalized.
This project demonstrated that a voice-command interface built on commodity offline recognition
hardware can control a quadrotor drone through the full mission cycle of takeoff, directional
flight, and landing without a conventional RC transmitter.
The ESP32 and VC02 combination achieved the required command latency budget, the ELRS
link maintained reliable packet delivery at 50~Hz across all flight sessions, and the SpeedyBee
F405~V5 platform proved stable on factory PID defaults without manual tuning.

\subsection{Summary of Accomplishments}

The following objectives from the high-level requirements were met:

\begin{itemize}
    \item A seven-word offline voice command vocabulary was implemented and mapped to
          flight-controller CRSF channel sequences with correct DVB-S2 framing and timing.
    \item The ELRS RF link was established using a RadioMaster RP2 reflashed as a CRSF
          transmitter, achieving automatic pairing via Binding Phrase without physical button
          binding.
    \item Stable hover was maintained for multiple seconds after each takeoff command, and
          directional pulse maneuvers produced observable displacements.
    \item Barometric altitude telemetry was decoded from the CRSF downlink and displayed on
          the OLED throughout flight, providing the ground operator with real-time situational
          awareness.
    \item % TODO: note landing success rate (e.g., X of Y landing attempts were successful).
\end{itemize}

\subsection{Ethical Considerations}

This project involved repeated outdoor and indoor flight of a motorized aerial vehicle with
rotating blades capable of causing injury.
The IEEE Code of Ethics \cite{ieee-ethics} requires engineers to hold the safety of the public
paramount and to be honest about the limitations of their work.
In adherence to these principles, all flight tests were conducted in restricted, supervised areas
with bystanders kept outside the rotor arc.
The voice control system incorporates a software interlock that rejects directional commands
unless the takeoff sequence has already been completed, preventing unintended ground-level
motor commands.
Future deployment in public spaces would require regulatory compliance with local UAV
operating rules and additional hardware failsafes such as a geofence or automatic return-to-home
function.

The project has broader relevance in the context of accessible human--machine interfaces.
Voice-driven UAV control reduces the physical skill barrier for operators with limited manual
dexterity and opens potential applications in disaster response and remote inspection where
operator hands may be otherwise engaged.
At the same time, increased accessibility of UAV control raises societal concerns around privacy
and airspace safety that future iterations of this work should address through both technical
design choices and engagement with regulatory frameworks.

\subsection{Future Work}

Several improvements would increase the practical utility of the system.

\textbf{Higher command update rate.}
The current 50~Hz CRSF output rate is sufficient for stable flight but is below the 250~Hz or
500~Hz rates achievable with ELRS.
Increasing the rate would reduce the latency between a voice command being received and the
first corresponding CRSF frame reaching the aircraft.

\textbf{Continuous control modality.}
Replacing discrete pulse commands with proportional control---for example, by mapping the
operator's head orientation or a wrist-worn IMU to roll and pitch channels in real time---would
enable more precise maneuvering and altitude hold.

\textbf{Safety failsafes.}
A hardware kill-switch input to the ESP32 and a CRSF failsafe channel configuration on the
flight controller would ensure the motors disarm if the voice link or USB power to the ground
unit is interrupted unexpectedly.

\textbf{Outdoor range characterization.}
All verification flights were conducted at short range.
A systematic link-budget test at increasing distances would characterize the operational envelope
of the RP2-based transmitter at the ISM\_2400 regulatory domain power level.
